\documentclass[11pt,a4paper]{article}
\usepackage{dsmorgan}

\title{Three Tangents to a Parabola}
\subtitle{A ratio identity, for every parabola and every third tangent}
\author{Drew M. Schaffner}

\begin{document}
\maketitle

\section{Statement}

Let $T_1$ and $T_2$ be tangents to a parabola at distinct points $P_1$ and $P_2$,
meeting at $P$. A third tangent $T$, at a point other than $P_1$ and $P_2$, meets
$T_1$ at $Q_1$ and $T_2$ at $Q_2$. Since $Q_1$ lies on the line $PP_1$, there is a
unique real number $t_1$ with $Q_1 = P + t_1 (P_1 - P)$; define $t_2$ in the same way
from $Q_2 = P + t_2 (P_2 - P)$. Thus $\abs{t_1} = PQ_1 / PP_1$ and
$\abs{t_2} = PQ_2 / PP_2$.

\begin{theorem}[Three tangents]\label{thm:main}
  In this situation $t_1 + t_2 = 1$. If the point of contact of $T$ lies on the arc
  between $P_1$ and $P_2$, then $0 < t_1, t_2 < 1$, and hence
  \[
    \frac{PQ_1}{PP_1} + \frac{PQ_2}{PP_2} = 1.
  \]
\end{theorem}

\section{Setting}

\begin{definition}[Similarity]\label{def:similarity}
  A \emph{similarity} of the plane is a map $S(X) = \lambda R X + v$, where
  $\lambda > 0$, $R$ is a rotation or reflection, and $v \in \R^2$.
\end{definition}

A similarity is an affine bijection. It maps lines to lines, preserves tangency
between lines and smooth curves, and multiplies every length by $\lambda$. In
particular, it preserves the coefficients $t_1$ and $t_2$, because
$S(P + t (P_1 - P)) = S(P) + t \bigl( S(P_1) - S(P) \bigr)$.

\begin{lemma}[All parabolas are similar]\label{lem:normal-form}
  Every parabola is the image of $y = x^2$ under a similarity.
\end{lemma}

\begin{proof}
  A parabola with focal distance $p > 0$ can be moved by a rotation and a
  translation to $y = x^2 / (4p)$. The scaling $(x, y) \mapsto (4p\,x, 4p\,y)$ takes
  $y = x^2$ to $y/(4p) = (x/(4p))^2$, that is, to $y = x^2/(4p)$.
\end{proof}

By \cref{lem:normal-form} and the invariance of $t_1$ and $t_2$, it suffices to
prove \cref{thm:main} for $y = x^2$.

\section{Idea of the proof}

On $y = x^2$ the tangents at $s$ and $u$ meet at $\bigl(\frac{s+u}{2}, su\bigr)$.
Measured along a tangent line, the parameters $t_1$ and $t_2$ are ratios of
$x$-coordinates, and they add to one by a single cancellation.

\section{Proof}

\begin{lemma}[Tangent intersections]\label{lem:meet}
  On $y = x^2$, the tangent at $(s, s^2)$ is $y = 2sx - s^2$. For $s \neq u$, the
  tangents at $s$ and $u$ meet at $\bigl(\frac{s+u}{2}, su\bigr)$.
\end{lemma}

\begin{proof}
  The slope at $x = s$ is $2s$, so the tangent is $y - s^2 = 2s(x - s)$. Two such
  lines meet where $2(s - u)x = s^2 - u^2$, that is, at $x = \frac{s+u}{2}$, and then
  $y = 2s \cdot \frac{s+u}{2} - s^2 = su$.
\end{proof}

\begin{proof}[Proof of \cref{thm:main}]
  Let the three points of contact have $x$-coordinates $a$, $b$, $c$, all distinct,
  with $P_1 = (a, a^2)$ and $P_2 = (b, b^2)$. By \cref{lem:meet},
  \[
    P = \Bigl(\frac{a+b}{2}, ab\Bigr), \qquad
    Q_1 = \Bigl(\frac{a+c}{2}, ac\Bigr), \qquad
    Q_2 = \Bigl(\frac{b+c}{2}, bc\Bigr).
  \]
  The line $T_1$ is not vertical, so $t_1$ is determined by $x$-coordinates alone:
  \[
    t_1 = \frac{x_{Q_1} - x_P}{x_{P_1} - x_P}
        = \frac{\frac{a+c}{2} - \frac{a+b}{2}}{a - \frac{a+b}{2}}
        = \frac{c - b}{a - b}.
  \]
  In the same way, on $T_2$,
  \[
    t_2 = \frac{\frac{b+c}{2} - \frac{a+b}{2}}{b - \frac{a+b}{2}} = \frac{c - a}{b - a}.
  \]
  Hence
  \[
    t_1 + t_2 = \frac{c - b}{a - b} - \frac{c - a}{a - b}
    = \frac{a - b}{a - b} = 1.
  \]
  If $c$ lies strictly between $a$ and $b$, then $c - b$ has the sign of $a - b$ and
  $\abs{c - b} < \abs{a - b}$, so $0 < t_1 < 1$; likewise $0 < t_2 < 1$. Then
  $PQ_1/PP_1 + PQ_2/PP_2 = t_1 + t_2 = 1$.
\end{proof}

\section{Remarks}

\begin{remark}
  When the third point of contact lies outside the arc $P_1P_2$, exactly one of
  $t_1$, $t_2$ is negative, and the identity becomes
  $\abs*{\frac{PQ_1}{PP_1} - \frac{PQ_2}{PP_2}} = 1$. The signed form $t_1 + t_2 = 1$
  covers both cases.
\end{remark}

\begin{remark}
  Since $t_1 = \frac{c - b}{a - b}$ depends on $c$ linearly, the map $Q_1 \mapsto Q_2$
  from $T_1$ to $T_2$ is affine. A moving tangent therefore cuts two fixed tangents
  in proportional segments. This is the classical construction of a parabola as the
  envelope of lines joining corresponding points of two proportionally divided
  segments, as in string art.
\end{remark}

\end{document}
